% Focused runtime contract for all public theorem insertion helpers.
% Build from the thesis/ directory after the template has initialized formats.
\input{./template/configure}

\begin{document}
\newcommand{\TypeoutTheoremDefault}[1]{%
  \typeout{NCKU-TEST-DEFAULT-#1: \csname GetTheorem#1FormatEnvironmentName\endcsname/\csname GetTheorem#1FormatShowText\endcsname/\csname GetTheorem#1FormatFollowCounter\endcsname}%
}
\TypeoutTheoremDefault{Definition}
\TypeoutTheoremDefault{Condition}
\TypeoutTheoremDefault{Problem}
\TypeoutTheoremDefault{Example}
\TypeoutTheoremDefault{Theorem}
\TypeoutTheoremDefault{Lemma}
\TypeoutTheoremDefault{Corollary}
\TypeoutTheoremDefault{Proposition}
\TypeoutTheoremDefault{Conjecture}
\TypeoutTheoremDefault{Proof}
\TypeoutTheoremDefault{Note}
\TypeoutTheoremDefault{Annotation}
\TypeoutTheoremDefault{Claim}
\TypeoutTheoremDefault{Case}
\TypeoutTheoremDefault{Acknowledgment}
\TypeoutTheoremDefault{Conclusion}
\TypeoutTheoremDefault{Criterion}
\TypeoutTheoremDefault{Assertion}
\TypeoutTheoremDefault{Question}
\TypeoutTheoremDefault{Hypothesis}
\TypeoutTheoremDefault{Summary}

% Freeze the two-key theorem-content parser independently of rendering.
\begingroup
\NCKUPrivateSetInsertTheoremKeys{\empty}
\typeout{NCKU-THEOREM-KEY-DEFAULTS: \TmpValueTitle/\TmpValueLabel}
\def\NCKUTheoremTitle{Expanded Theorem Title}
\def\NCKUTheoremLabel{ncku:theorem-expanded}
\NCKUPrivateSetInsertTheoremKeys{%
  title=\NCKUTheoremTitle,
  label=\NCKUTheoremLabel}
\typeout{NCKU-THEOREM-KEY-EXPANDED: \TmpValueTitle/\TmpValueLabel}
\NCKUPrivateSetInsertTheoremKeys{\empty}
\typeout{NCKU-THEOREM-KEY-RESET: \TmpValueTitle/\TmpValueLabel}
\endgroup

\stepcounter{section}

\InsertDefinition[label={ncku:test:definition}]{NCKU Definition Body.}
\InsertCondition[label={ncku:test:condition}]{NCKU Condition Body.}
\InsertProblem[label={ncku:test:problem}]{NCKU Problem Body.}
\InsertExample[label={ncku:test:example}]{NCKU Example Body.}
\InsertTheorem[title={Named Theorem},label={ncku:test:theorem-first}]{NCKU Theorem Body One.}
\InsertLemma[label={ncku:test:lemma}]{NCKU Lemma Body.}
\InsertCorollary[label={ncku:test:corollary}]{NCKU Corollary Body.}
\InsertProposition[label={ncku:test:proposition}]{NCKU Proposition Body.}
\InsertConjecture[label={ncku:test:conjecture}]{NCKU Conjecture Body.}
\InsertCriterion[label={ncku:test:criterion}]{NCKU Criterion Body.}
\InsertAssertion[label={ncku:test:assertion}]{NCKU Assertion Body.}
\InsertQuestion[label={ncku:test:question}]{NCKU Question Body.}
\InsertHypothesis[label={ncku:test:hypothesis}]{NCKU Hypothesis Body.}
\InsertTheorem[label={ncku:test:theorem-second}]{NCKU Theorem Body Two.}

\InsertProof{NCKU Proof Body.}
\InsertNote{NCKU Note Body.}
\InsertAnnotation{NCKU Annotation Body.}
\InsertClaim{NCKU Claim Body.}
\InsertCase{NCKU Case Body.}
\InsertAcknowledgment{NCKU Acknowledgment Body.}
\InsertConclusion{NCKU Conclusion Body.}
\InsertSummary{NCKU Summary Body.}

\stepcounter{section}
\InsertTheorem[label={ncku:test:theorem-reset}]{NCKU Theorem Reset Body.}

References: \ref{ncku:test:definition}, \ref{ncku:test:theorem-first}, \ref{ncku:test:theorem-second}, \ref{ncku:test:theorem-reset}.
Named reference: \nameref{ncku:test:theorem-first}.

% Exercise every public setter route after rendering so getter mutation cannot
% affect the established PDF contract above.
\newcommand{\ExerciseTheoremFormatRoute}[1]{%
  \SetTheoremFormat[#1]{ShowText = {Registry#1}}%
  \typeout{NCKU-TEST-ROUTE-#1: \csname GetTheorem#1FormatShowText\endcsname}%
}
\ExerciseTheoremFormatRoute{Definition}
\ExerciseTheoremFormatRoute{Condition}
\ExerciseTheoremFormatRoute{Problem}
\ExerciseTheoremFormatRoute{Example}
\ExerciseTheoremFormatRoute{Theorem}
\ExerciseTheoremFormatRoute{Lemma}
\ExerciseTheoremFormatRoute{Corollary}
\ExerciseTheoremFormatRoute{Proposition}
\ExerciseTheoremFormatRoute{Conjecture}
\ExerciseTheoremFormatRoute{Proof}
\ExerciseTheoremFormatRoute{Note}
\ExerciseTheoremFormatRoute{Annotation}
\ExerciseTheoremFormatRoute{Claim}
\ExerciseTheoremFormatRoute{Case}
\ExerciseTheoremFormatRoute{Acknowledgment}
\ExerciseTheoremFormatRoute{Conclusion}
\ExerciseTheoremFormatRoute{Criterion}
\ExerciseTheoremFormatRoute{Assertion}
\ExerciseTheoremFormatRoute{Question}
\ExerciseTheoremFormatRoute{Hypothesis}
\ExerciseTheoremFormatRoute{Summary}

\SetTheoremFormat[Unknown]{ShowText = {UnexpectedUnknownRoute}}
\ifcsname GetTheoremUnknownFormatShowText\endcsname
  \typeout{NCKU-TEST-UNKNOWN-ROUTE: FAIL}
\else
  \typeout{NCKU-TEST-UNKNOWN-ROUTE: no-op}
\fi
\SetTheoremFormat[]{ShowText = {UnexpectedEmptyRoute}}
\ifcsname GetTheoremFormatShowText\endcsname
  \typeout{NCKU-TEST-EMPTY-ROUTE: FAIL}
\else
  \typeout{NCKU-TEST-EMPTY-ROUTE: no-op}
\fi

% Exercise every registered counter target after the setter routes reset raw
% defaults. Known targets resolve to an effective counter; unknown/empty targets
% preserve the source value. The self-route must freeze rather than alias itself.
\newcommand{\ExerciseTheoremCounterRoute}[2]{%
  \SetTheoremFormat[#1]{FollowCounter = Section}%
  \MappingTheoremCounter[#1]{#2}%
  \typeout{NCKU-TEST-COUNTER-ROUTE-#2: #1/\csname GetTheorem#1FormatFollowCounter\endcsname}%
}
\ExerciseTheoremCounterRoute{Definition}{Definition}
\ExerciseTheoremCounterRoute{Definition}{Condition}
\ExerciseTheoremCounterRoute{Definition}{Problem}
\ExerciseTheoremCounterRoute{Definition}{Example}
\ExerciseTheoremCounterRoute{Definition}{Theorem}
\ExerciseTheoremCounterRoute{Definition}{Lemma}
\ExerciseTheoremCounterRoute{Definition}{Corollary}
\ExerciseTheoremCounterRoute{Definition}{Proposition}
\ExerciseTheoremCounterRoute{Definition}{Conjecture}
\ExerciseTheoremCounterRoute{Definition}{Proof}
\ExerciseTheoremCounterRoute{Definition}{Note}
\ExerciseTheoremCounterRoute{Definition}{Annotation}
\ExerciseTheoremCounterRoute{Definition}{Claim}
\ExerciseTheoremCounterRoute{Definition}{Case}
\ExerciseTheoremCounterRoute{Definition}{Acknowledgment}
\ExerciseTheoremCounterRoute{Definition}{Conclusion}
\ExerciseTheoremCounterRoute{Definition}{Criterion}
\ExerciseTheoremCounterRoute{Definition}{Assertion}
\ExerciseTheoremCounterRoute{Definition}{Question}
\ExerciseTheoremCounterRoute{Definition}{Hypothesis}
\ExerciseTheoremCounterRoute{Definition}{Summary}
\ExerciseTheoremCounterRoute{Definition}{Section}

\SetTheoremFormat[Definition]{FollowCounter = Section}
\MappingTheoremCounter[Definition]{Unknown}
\typeout{NCKU-TEST-COUNTER-UNKNOWN: \GetTheoremDefinitionFormatFollowCounter}
\SetTheoremFormat[Definition]{FollowCounter = Section}
\MappingTheoremCounter[Definition]{}
\typeout{NCKU-TEST-COUNTER-EMPTY: \GetTheoremDefinitionFormatFollowCounter}

\SetTheoremFormat[Definition]{FollowCounter = \empty}
\SetTheoremFormat[Theorem]{FollowCounter = Definition}
\SetTheoremFormat[Problem]{FollowCounter = Theorem}
\MappingTheoremCounter[Problem]{Theorem}
\typeout{NCKU-TEST-COUNTER-CHAIN: \GetTheoremProblemFormatFollowCounter}

\typeout{NCKU-TEST-PASS: all 21 public theorem insertion helpers compiled}
\end{document}
